\chapter{1936:Peter Debye}

\label{chap:chemistry1936}

The Nobel Prize in Chemistry for the year 1936 was awarded to Peter Debye.

\textbf{Motivation:}

for his contributions to our knowledge of molecular structure through his investigations on dipole moments and on the diffraction of X-rays and electrons in gases

\section{Peter Debye}

\subsection*{Early Life and Education}

Peter Joseph William Debye was born on 1884-03-24 in Maastricht, Netherlands.
He was a Dutch-American physical chemist awarded the Nobel Prize in Chemistry in 1936 for his contributions to the understanding of molecular structure through his investigations on dipole moments and on the diffraction of X-rays and electrons in gases.

\subsection*{Career and Major Research}

Debye studied at the Technische Hochschule in Aachen and at the University of Munich, where he received his doctorate in 1908 under Arnold Sommerfeld.
He held professorships in Zürich in 1911, Utrecht in 1912, Göttingen in 1914, Zürich in 1920, and Leipzig in 1927.
In 1934 he succeeded Albert Einstein as director of the Kaiser Wilhelm Institute for Physics in Berlin.
In 1940 he moved to Cornell University in the United States, where he was professor of chemistry until 1950, and he became an American citizen in 1946.
He contributed to the theory of specific heat, developed with Erich Hückel the theory of strong electrolytes, and devised with Paul Scherrer the powder method of X-ray diffraction.
He developed the use of electric dipole moments as a means of determining molecular structure, and he applied both X-ray and electron diffraction to gases in order to measure the distances between atoms in molecules.

\subsection*{Nobel Prize-winning Work}

The work for which Peter Debye was awarded the Nobel Prize in Chemistry in 1936 was described by the Nobel Committee as: "for his contributions to our knowledge of molecular structure through his investigations on dipole moments and on the diffraction of X-rays and electrons in gases".

By measuring the dipole moments of molecules in solution, Debye obtained information about the arrangement of atoms within the molecule, and by diffracting X-rays and electrons from gases he measured the distances between the atoms in simple molecules.

\subsection*{Significance and Impact}

Debye's methods gave chemists quantitative information about the geometry of molecules, and the measurement of dipole moments became a standard tool of physical chemistry.
The unit of the electric dipole moment, the debye (D), is named in his honour.

\subsection*{Later Career and Honors}

Debye remained at Cornell University until his death, serving as head of the chemistry department.
He received the Lorentz Medal, the Franklin Medal, and many honorary doctorates.
He died on 1966-11-02 in Ithaca, New York, United States.

\subsection*{Legacy}

The legacy of Peter Debye endures through the lasting impact of his scientific contributions.
His theories and methods continue to be used throughout physical chemistry and molecular physics.

\subsection*{Modern Developments}

Debye's model of the dipole and his theory of solutions and dielectrics remain central to modern physical chemistry and materials science. Debye–Hückel theory governs electrolyte solutions, while Debye scattering and the Debye temperature are used in the analysis of crystals and heat capacities. The Debye–Scherrer powder method is a standard tool in X-ray and neutron diffraction.

\subsection*{Selected Publications}

\begin{itemize}

\item Debye, P., \& Scherrer, P. (1916). "Interferenzen an regellos orientierten Teilchen im Röntgenlicht." Physikalische Zeitschrift.

\item Debye, P., \& Hückel, E. (1923). "Zur Theorie der Elektrolyte." Physikalische Zeitschrift.

\item Debye, P. (1929). "Polar Molecules." New York: Chemical Catalog Company.

\end{itemize}

\clearpage

\section*{Chapter Summary}

The 1936 prize recognized Peter Debye for his contributions to the knowledge of molecular structure through his investigations on dipole moments and on the diffraction of X-rays and electrons in gases.
